\scanpage{69}
% archive foliation: 68

\begin{center}- 30 -\end{center}

D'autre part, on a, en tenant compte de (26)
\[
\begin{aligned}
\frac d{dt}\bigl(tU_{kj}(\mathbf{x}_0+t(\mathbf{x}-\mathbf{x}_0);\mathbf{v}(t))\bigr)
={}&U_{kj}(\mathbf{x}_0+t(\mathbf{x}-\mathbf{x}_0);\mathbf{v}(t))+\\
&+t\sum_i(x_i-x_{i0})\frac{\partial U_{kj}}{\partial x_i}
(\mathbf{x}_0+t(\mathbf{x}-\mathbf{x}_0);\mathbf{v}(t))\\
&+t\sum_{i,h}(x_i-x_{i0})\frac{\partial U_{kj}}{\partial y_h}
(\mathbf{x}_0+t(\mathbf{x}-\mathbf{x}_0);\mathbf{v}(t))U_{ih}(\mathbf{x}_0+t(\mathbf{x}-\mathbf{x}_0);\mathbf{v}(t))
\end{aligned}
\]
d'où, en faisant intervenir (24), on voit que les fonctions
\[
w_{kj}(t)=\frac{\partial v_j}{\partial x_k}
-tU_{kj}(\mathbf{x}_0+t(\mathbf{x}-\mathbf{x}_0);\mathbf{v}(t))
\qquad(1\leq j\leq q,\ 1\leq k\leq p)
\]
\mathbf{v}érifient le système linéaire et homogène des $pq$ équations différentielles
\[
(27)\qquad
\frac{dw_{kj}}{dt}=\sum_{i,h}(x_i-x_{i0})\frac{\partial U_{ij}}{\partial y_h}w_{kh}
\]
Or $w_{kj}(0)=0$, car pour $t=0$, $\mathbf{v}(0;\mathbf{x}_0,\mathbf{x},\mathbf{y}_0)=\mathbf{y}_0$, donc $\frac{\partial v_j}{\partial x_k}(0)=0$ ;
comme $0$ est la seule solution de (27) qui s'annule pour $t=0$, on a identiquement $w_{kj}(t)=0$ pour $|t|<2$, et en particulier pour $t=1$ ;
mais si on pose $\mathbf{u}(\mathbf{x},\mathbf{x}_0,\mathbf{y}_0)=\mathbf{v}(1;\mathbf{x}_0,\mathbf{x},\mathbf{y}_0)$, cela signifie que l'on a
\[
\frac{\partial u_j}{\partial x_i}=U_{ij}(\mathbf{x};\mathbf{u}(\mathbf{x}))
\qquad(1\leq i\leq p,\ 1\leq j\leq q)
\]
pour $\mathbf{x}\in U$. De ces relations on déduit que si $\mathbf{z}(t)=\mathbf{u}(\mathbf{x}_0+t(\mathbf{x}-\mathbf{x}_0))$, on a
\[
\mathbf{z}'(t)=\uline{U}(\mathbf{x}_0+t(\mathbf{x}-\mathbf{x}_0);\mathbf{z}(t)).(\mathbf{x}-\mathbf{x}_0)
\]
pour $|t|<2$, et $\mathbf{z}(1)=\mathbf{v}(1)$ ; donc $\mathbf{z}(t)=\mathbf{v}(t)$ pour $|t|\leq1$, et en particulier $\mathbf{u}(\mathbf{x}_0)=\mathbf{z}(0)=\mathbf{v}(0)=\mathbf{y}_0$ ; ce qui achève la démonstration.

On pourrait évidemment ici passer de la forme non homogène à la forme homogène du th. de Frobenius, avec l'algèbre extérieure, parler de la primitive (locale) d'une forme extérieure holomorphe, etc..; est-ce le lieu, ou,vaut-il mieux attendre les variétés analytiques pour le faire ?
